\documentclass[11pt]{article}
\usepackage[letterpaper,margin=0.72in]{geometry}
\usepackage[T1]{fontenc}
\usepackage[utf8]{inputenc}
\usepackage{lmodern}
\usepackage{microtype}
\usepackage{xcolor}
\usepackage{graphicx}
\usepackage{booktabs}
\usepackage{longtable}
\usepackage{tabularx}
\usepackage{array}
\usepackage{enumitem}
\usepackage[most]{tcolorbox}
\usepackage{fancyhdr}
\usepackage{titlesec}
\usepackage{hyperref}
\usepackage{xurl}
\usepackage{lastpage}
\usepackage{ragged2e}
\usepackage{multicol}
\usepackage{amssymb}
\definecolor{BAOrange}{HTML}{E87521}
\definecolor{BANavy}{HTML}{17324D}
\definecolor{BASlate}{HTML}{5B6573}
\definecolor{BALight}{HTML}{F4F6F8}
\definecolor{BAGold}{HTML}{D7A62A}
\hypersetup{colorlinks=true,linkcolor=BANavy,urlcolor=BAOrange,citecolor=BANavy,pdftitle={Cornell-Style Chapter Notes}}
\setlength{\parindent}{0pt}
\setlength{\parskip}{4pt}
\setlist{nosep,leftmargin=*}
\renewcommand{\arraystretch}{1.18}
\newcolumntype{Y}{>{\RaggedRight\arraybackslash}X}
\newcolumntype{L}[1]{>{\RaggedRight\arraybackslash}p{#1}}
\pagestyle{fancy}
\setlength{\headheight}{14pt}
\fancyhf{}
\lhead{\textcolor{BANavy}{\small BABOK v3 Study Series}}
\rhead{\textcolor{BASlate}{\small Cornell-Style Chapter Notes}}
\cfoot{\small \thepage\ / \pageref{LastPage}}
\renewcommand{\headrulewidth}{0.4pt}
\titleformat{\section}{\Large\bfseries\color{BANavy}}{\thesection}{0.7em}{}
\titleformat{\subsection}{\large\bfseries\color{BAOrange}}{\thesubsection}{0.6em}{}
\titleformat{\subsubsection}{\normalsize\bfseries\color{BANavy}}{\thesubsubsection}{0.5em}{}
\newtcolorbox{keybox}[1][]{enhanced,breakable,colback=BALight,colframe=BANavy,boxrule=0.7pt,arc=1.5mm,left=2mm,right=2mm,top=1.5mm,bottom=1.5mm,#1}
\newtcolorbox{practicebox}[1][]{enhanced,breakable,colback=white,colframe=BAOrange,boxrule=0.8pt,arc=1.5mm,left=2mm,right=2mm,top=1.5mm,bottom=1.5mm,#1}
\newtcolorbox{warningbox}[1][]{enhanced,breakable,colback=white,colframe=BAGold,boxrule=0.8pt,arc=1.5mm,left=2mm,right=2mm,top=1.5mm,bottom=1.5mm,#1}
\begin{document}
\begin{titlepage}
\centering
\vspace*{1.15in}
{\Huge\bfseries\color{BANavy} Cornell-Style Chapter Notes\par}
\vspace{0.35in}
{\Large\color{BAOrange} Eleven-chapter review and retention guide\par}
\vspace{0.55in}
\rule{0.78\textwidth}{1.2pt}\par
\vspace{0.35in}
{\large Technical reference and study artifact based on the uploaded \emph{A Guide to the Business Analysis Body of Knowledge (BABOK Guide), Version 3}.\par}
\vfill
\begin{keybox}
\textbf{Source discipline.} BABOK terminology and organization are preserved where practical. Explanations are paraphrased for study and engineering use. Any engineering, security, Agile, UML/BPMN, or enterprise-architecture crosswalks are practitioner mappings, not normative BABOK statements.
\end{keybox}
\vspace{0.35in}
{\small Prepared October 2026\par}
\end{titlepage}
\tableofcontents
\newpage
\section{How to Use These Notes}
Each chapter is organized in Cornell form: cues and questions on the left, study notes on the right, followed by a synthesis summary. The notes intentionally paraphrase rather than reproduce the source text. Use the cue column for active recall: cover the notes, answer each cue from memory, then compare.
\section{Chapter 1: Introduction}
\begin{longtable}{L{0.26\textwidth} L{0.64\textwidth}}
\toprule
\textbf{Cue / Question} & \textbf{Notes}\\
\midrule
\endhead
What is the guide for? & A common framework for discussing and performing business analysis. It defines widely accepted practices, but it is not a mandated sequence or a one-size-fits-all process.\\
What is business analysis? & Enabling change by defining needs and recommending solutions that deliver stakeholder value. Analysis can occur at strategic, tactical, operational, project, and continuous-improvement levels.\\
Who performs BA? & Anyone performing business analysis tasks, regardless of title. Product managers, product owners, systems analysts, architects, consultants, and requirements engineers can all perform BA work.\\
How is the guide structured? & Core concepts provide the conceptual foundation; six knowledge areas group tasks; competencies describe enabling skills; techniques provide methods; perspectives adapt the work to context.\\
Are tasks sequential? & No. Tasks may be performed sequentially, iteratively, or concurrently when needed inputs are available.\\
Task anatomy & Each task is described through purpose, description, inputs, elements, guidelines/tools, techniques, stakeholders, and outputs.\\
\bottomrule
\end{longtable}
\begin{keybox}\textbf{Chapter summary.} Treat the guide as a flexible body of practice. Select and tailor tasks and techniques to the initiative rather than forcing a lifecycle onto the work.\end{keybox}
\subsection*{Active-recall prompts}
\begin{enumerate}
\item Explain the central purpose of Chapter 1 without using the chapter title.
\item Name the most important decision or failure this chapter helps prevent.
\item Give one example of how the chapter would apply to a software product change.
\item Identify one artifact or model that could provide evidence that the work was performed well.
\end{enumerate}
\section{Chapter 2: Business Analysis Key Concepts}
\begin{longtable}{L{0.26\textwidth} L{0.64\textwidth}}
\toprule
\textbf{Cue / Question} & \textbf{Notes}\\
\midrule
\endhead
BACCM & The Business Analysis Core Concept Model connects six equally important concepts: Change, Need, Solution, Stakeholder, Value, and Context.\\
Change & A transformation undertaken in response to a need. Changes are analyzed in terms of the value they seek and the context in which they occur.\\
Need & A problem or opportunity. Needs can motivate change, while changes can also create new needs.\\
Solution & A particular way of satisfying one or more needs within a context.\\
Stakeholder & A person or group related to the need, change, or solution; analysis considers interest, impact, and influence.\\
Value & The worth, usefulness, gain, or loss experienced by stakeholders. Value may be tangible or intangible and is often comparative.\\
Context & The relevant circumstances surrounding the change: culture, technology, policy, processes, markets, constraints, and other environmental factors.\\
Requirement vs design & A requirement represents a need; a design represents a solution. The distinction depends on how the information is used and by whom.\\
Requirement classes & Business, stakeholder, solution (functional and non-functional), and transition requirements.\\
\bottomrule
\end{longtable}
\begin{keybox}\textbf{Chapter summary.} Use the six core concepts as a completeness check: every significant change should be explainable in terms of need, stakeholders, context, solution, and value.\end{keybox}
\subsection*{Active-recall prompts}
\begin{enumerate}
\item Explain the central purpose of Chapter 2 without using the chapter title.
\item Name the most important decision or failure this chapter helps prevent.
\item Give one example of how the chapter would apply to a software product change.
\item Identify one artifact or model that could provide evidence that the work was performed well.
\end{enumerate}
\section{Chapter 3: Business Analysis Planning and Monitoring}
\begin{longtable}{L{0.26\textwidth} L{0.64\textwidth}}
\toprule
\textbf{Cue / Question} & \textbf{Notes}\\
\midrule
\endhead
Approach & Select and tailor the way analysis will be conducted, including formality, timing, techniques, deliverables, and relationship to the delivery method.\\
Stakeholder engagement & Identify relevant stakeholders, their needs from analysis, what analysis needs from them, and how collaboration will work.\\
Governance & Define who can decide, review, prioritize, approve, and authorize changes to requirements and designs.\\
Information management & Decide how BA information will be organized, stored, accessed, versioned, traced, and retained.\\
Performance improvement & Measure the quality and effectiveness of BA work; use findings and lessons to improve future performance.\\
\bottomrule
\end{longtable}
\begin{keybox}\textbf{Chapter summary.} Planning is not merely project administration. It establishes the decision system, information system, stakeholder model, and improvement loop that make later analysis reliable.\end{keybox}
\subsection*{Active-recall prompts}
\begin{enumerate}
\item Explain the central purpose of Chapter 3 without using the chapter title.
\item Name the most important decision or failure this chapter helps prevent.
\item Give one example of how the chapter would apply to a software product change.
\item Identify one artifact or model that could provide evidence that the work was performed well.
\end{enumerate}
\section{Chapter 4: Elicitation and Collaboration}
\begin{longtable}{L{0.26\textwidth} L{0.64\textwidth}}
\toprule
\textbf{Cue / Question} & \textbf{Notes}\\
\midrule
\endhead
Elicitation is ongoing & Elicitation is not a single requirements phase. It continues while analysis work continues and may be planned or just-in-time.\\
Prepare & Clarify objectives and scope, select techniques, identify participants, and prepare supporting information.\\
Conduct & Use conversations, research, observation, experiments, workshops, or other techniques to obtain information relevant to the change.\\
Confirm & Reconcile elicitation results for accuracy and consistency before treating them as dependable business analysis information.\\
Communicate & Tailor the form, level of detail, and medium to the audience so stakeholders share an understanding.\\
Collaborate & Manage stakeholder participation, trust, conflict, and commitment toward common goals.\\
\bottomrule
\end{longtable}
\begin{keybox}\textbf{Chapter summary.} High-quality requirements depend on high-quality collaboration. Confirmation and communication are as important as the initial act of gathering information.\end{keybox}
\subsection*{Active-recall prompts}
\begin{enumerate}
\item Explain the central purpose of Chapter 4 without using the chapter title.
\item Name the most important decision or failure this chapter helps prevent.
\item Give one example of how the chapter would apply to a software product change.
\item Identify one artifact or model that could provide evidence that the work was performed well.
\end{enumerate}
\section{Chapter 5: Requirements Life Cycle Management}
\begin{longtable}{L{0.26\textwidth} L{0.64\textwidth}}
\toprule
\textbf{Cue / Question} & \textbf{Notes}\\
\midrule
\endhead
Life cycle & Requirements and designs are managed from initial representation of a need through development, operation, change, and eventual retirement.\\
Trace & Maintain backward, forward, and peer relationships so impact, coverage, allocation, gaps, and implementation status can be understood.\\
Maintain & Keep requirements accurate, current, consistent, accessible, and reusable where appropriate.\\
Prioritize & Use value, urgency, risk, dependencies, and other agreed factors to order analysis and delivery work.\\
Assess changes & Analyze the impact of proposed requirement or design changes before accepting or rejecting them.\\
Approve & Obtain appropriate agreement and authorization according to the governance model.\\
\bottomrule
\end{longtable}
\begin{keybox}\textbf{Chapter summary.} Life-cycle management converts requirements from static statements into governed, traceable assets that can survive change.\end{keybox}
\subsection*{Active-recall prompts}
\begin{enumerate}
\item Explain the central purpose of Chapter 5 without using the chapter title.
\item Name the most important decision or failure this chapter helps prevent.
\item Give one example of how the chapter would apply to a software product change.
\item Identify one artifact or model that could provide evidence that the work was performed well.
\end{enumerate}
\section{Chapter 6: Strategy Analysis}
\begin{longtable}{L{0.26\textwidth} L{0.64\textwidth}}
\toprule
\textbf{Cue / Question} & \textbf{Notes}\\
\midrule
\endhead
Current state & Understand the business need and the existing capabilities, processes, structures, technology, policies, constraints, and performance that shape the change.\\
Future state & Describe measurable conditions, goals, objectives, capabilities, and outcomes that would satisfy the need.\\
Risk & Identify uncertainty affecting transition or future-state value and determine how it should influence decisions.\\
Change strategy & Compare approaches for closing the current-to-future-state gap, including transition states, dependencies, costs, risks, and value.\\
Continuous analysis & Strategy analysis can recur as context, needs, assumptions, or evidence changes.\\
\bottomrule
\end{longtable}
\begin{keybox}\textbf{Chapter summary.} Strategy analysis connects a business need to a credible transition path. It prevents requirements work from becoming detached from enterprise outcomes.\end{keybox}
\subsection*{Active-recall prompts}
\begin{enumerate}
\item Explain the central purpose of Chapter 6 without using the chapter title.
\item Name the most important decision or failure this chapter helps prevent.
\item Give one example of how the chapter would apply to a software product change.
\item Identify one artifact or model that could provide evidence that the work was performed well.
\end{enumerate}
\section{Chapter 7: Requirements Analysis and Design Definition}
\begin{longtable}{L{0.26\textwidth} L{0.64\textwidth}}
\toprule
\textbf{Cue / Question} & \textbf{Notes}\\
\midrule
\endhead
Specify/model & Turn elicitation results into coherent requirements and design representations appropriate to the audience and problem.\\
Verify & Assess quality: clarity, consistency, completeness, usability, feasibility of analysis, and fitness for intended use.\\
Validate & Assess value alignment: does this requirement or design support the business need and desired outcomes?\\
Requirements architecture & Organize requirements and relationships so the complete set supports the objectives and can be navigated and reasoned about.\\
Design options & Develop alternative solution approaches and allocate requirements across components.\\
Recommend solution & Compare potential value and other considerations to identify the most appropriate solution option.\\
\bottomrule
\end{longtable}
\begin{keybox}\textbf{Chapter summary.} This knowledge area is where raw information becomes structured, testable, value-aligned requirements and design options.\end{keybox}
\subsection*{Active-recall prompts}
\begin{enumerate}
\item Explain the central purpose of Chapter 7 without using the chapter title.
\item Name the most important decision or failure this chapter helps prevent.
\item Give one example of how the chapter would apply to a software product change.
\item Identify one artifact or model that could provide evidence that the work was performed well.
\end{enumerate}
\section{Chapter 8: Solution Evaluation}
\begin{longtable}{L{0.26\textwidth} L{0.64\textwidth}}
\toprule
\textbf{Cue / Question} & \textbf{Notes}\\
\midrule
\endhead
Actual solution exists & Unlike purely prospective analysis, solution evaluation examines something already operating in some form: prototype, pilot, beta, component, or production release.\\
Measure & Choose measures aligned with goals and collect reliable performance data.\\
Analyze & Interpret performance measures to determine whether the solution is delivering expected value.\\
Solution limitations & Identify constraints inside the solution that reduce value.\\
Enterprise limitations & Identify organizational, process, capability, policy, or environmental constraints outside the solution.\\
Increase value & Recommend changes, workarounds, optimization, replacement, or retirement actions that improve realized value.\\
\bottomrule
\end{longtable}
\begin{keybox}\textbf{Chapter summary.} Evaluation closes the loop between expected and realized value. A solution is successful only when evidence shows that stakeholder and enterprise outcomes are being achieved.\end{keybox}
\subsection*{Active-recall prompts}
\begin{enumerate}
\item Explain the central purpose of Chapter 8 without using the chapter title.
\item Name the most important decision or failure this chapter helps prevent.
\item Give one example of how the chapter would apply to a software product change.
\item Identify one artifact or model that could provide evidence that the work was performed well.
\end{enumerate}
\section{Chapter 9: Underlying Competencies}
\begin{longtable}{L{0.26\textwidth} L{0.64\textwidth}}
\toprule
\textbf{Cue / Question} & \textbf{Notes}\\
\midrule
\endhead
Analytical thinking & Creative thinking, decision making, learning, problem solving, systems thinking, and conceptual thinking support sound analysis.\\
Behavioural characteristics & Ethics, accountability, trustworthiness, adaptability, and personal organization affect credibility and effectiveness.\\
Business knowledge & Understanding business acumen, industry context, organization, solution knowledge, methodology, and stakeholder context improves judgment.\\
Communication skills & Verbal, non-verbal, written, and listening skills are central to elicitation, negotiation, documentation, and decision support.\\
Interaction skills & Facilitation, leadership, teamwork, negotiation, conflict resolution, and teaching help groups reach decisions.\\
Tools and technology & Productivity, requirements management, modelling, collaboration, and communication technologies amplify BA effectiveness when selected appropriately.\\
\bottomrule
\end{longtable}
\begin{keybox}\textbf{Chapter summary.} Competencies are enabling capabilities rather than a separate workflow. Weak communication, facilitation, or analytical skill can undermine otherwise correct techniques.\end{keybox}
\subsection*{Active-recall prompts}
\begin{enumerate}
\item Explain the central purpose of Chapter 9 without using the chapter title.
\item Name the most important decision or failure this chapter helps prevent.
\item Give one example of how the chapter would apply to a software product change.
\item Identify one artifact or model that could provide evidence that the work was performed well.
\end{enumerate}
\section{Chapter 10: Techniques}
\begin{longtable}{L{0.26\textwidth} L{0.64\textwidth}}
\toprule
\textbf{Cue / Question} & \textbf{Notes}\\
\midrule
\endhead
What techniques are & Methods used to perform business analysis tasks. The guide describes common techniques but does not present an exhaustive list.\\
Selection principle & Choose techniques based on context, stakeholders, uncertainty, information needs, cost, timing, and desired outputs.\\
Combination & Techniques can be used independently or in combination. A workshop may contain brainstorming, process modelling, prioritization, and review activities.\\
Modelling families & Process, data, state, sequence, scope, concept, and organizational models represent different views of a problem or solution.\\
Decision families & Decision analysis, decision modelling, financial analysis, risk analysis, business cases, and prioritization support choices and governance.\\
Elicitation families & Interviews, workshops, observation, focus groups, surveys, document analysis, collaborative games, and prototyping obtain or validate information.\\
\bottomrule
\end{longtable}
\begin{keybox}\textbf{Chapter summary.} Do not choose a technique because it is familiar. Choose it because it produces the type of evidence or shared understanding needed for the current task.\end{keybox}
\subsection*{Active-recall prompts}
\begin{enumerate}
\item Explain the central purpose of Chapter 10 without using the chapter title.
\item Name the most important decision or failure this chapter helps prevent.
\item Give one example of how the chapter would apply to a software product change.
\item Identify one artifact or model that could provide evidence that the work was performed well.
\end{enumerate}
\section{Chapter 11: Perspectives}
\begin{longtable}{L{0.26\textwidth} L{0.64\textwidth}}
\toprule
\textbf{Cue / Question} & \textbf{Notes}\\
\midrule
\endhead
Purpose of perspectives & Perspectives adapt business analysis to a context without replacing the underlying knowledge areas and tasks.\\
Agile & Emphasizes continuous analysis, incremental delivery, just-in-time detail, prioritized work, and close collaboration.\\
Business Intelligence & Focuses analysis on information, analytics, data quality, decision support, and value derived from data.\\
Information Technology & Focuses on the effects of change on IT systems and on translating business needs into technology-oriented solutions.\\
Business Architecture & Emphasizes strategy, capabilities, operating models, value, enterprise scope, and alignment across initiatives.\\
Business Process Management & Applies a process-centric lens to how work is performed and improved across organizational boundaries.\\
Multiple perspectives & Complex initiatives can use several perspectives at once; for example, an agile technology initiative that also changes processes and enterprise capabilities.\\
\bottomrule
\end{longtable}
\begin{keybox}\textbf{Chapter summary.} Perspectives are lenses, not silos. Use the combination that best explains the initiative and helps stakeholders reason about change.\end{keybox}
\subsection*{Active-recall prompts}
\begin{enumerate}
\item Explain the central purpose of Chapter 11 without using the chapter title.
\item Name the most important decision or failure this chapter helps prevent.
\item Give one example of how the chapter would apply to a software product change.
\item Identify one artifact or model that could provide evidence that the work was performed well.
\end{enumerate}
\section{Whole-Guide Synthesis}
The chapters form a connected reasoning system: concepts establish the vocabulary; planning establishes how decisions and information will be managed; elicitation supplies information; strategy frames the change and expected value; requirements life-cycle management preserves alignment over time; analysis and design definition turns information into coherent requirements and options; solution evaluation compares expected with realized value; competencies, techniques, and perspectives enable and tailor all of the above.

\begin{practicebox}
\textbf{Exam and practitioner memory hook:} \emph{Why change? Who is affected? What value is sought? What information do we need? How will we represent and govern it? Which solution best fits? What happened after implementation?}
\end{practicebox}

\section*{Primary Source}
\addcontentsline{toc}{section}{Primary Source}
International Institute of Business Analysis (IIBA), \emph{A Guide to the Business Analysis Body of Knowledge (BABOK Guide)}, Version 3, 2015. This document is a study and practitioner synthesis of the user-supplied copy; it is not an official IIBA publication.
\end{document}
